\documentclass[12pt,a4paper]{article}

% =============================================================================
% PAQUETES
% =============================================================================
\usepackage[utf8]{inputenc}
\usepackage[T1]{fontenc}
\usepackage[spanish]{babel}
\usepackage{geometry}
\usepackage{graphicx}
\usepackage{xcolor}
\usepackage{booktabs}
\usepackage{longtable}
\usepackage{array}
\usepackage{multirow}
\usepackage{float}
\usepackage{caption}
\usepackage{subcaption}
\usepackage{amsmath}
\usepackage{amssymb}
\usepackage{enumitem}
\usepackage{fancyhdr}
\usepackage{titlesec}
\usepackage{setspace}
\usepackage{tikz}
\usepackage{pgfplots}
\pgfplotsset{compat=1.18}

% =============================================================================
% CONFIGURACION DE PAGINA
% =============================================================================
\geometry{
    left=2.5cm,
    right=2.5cm,
    top=2.5cm,
    bottom=2.5cm
}

% =============================================================================
% COLORES
% =============================================================================
\definecolor{azuloscuro}{RGB}{0, 51, 102}
\definecolor{azulclaro}{RGB}{70, 130, 180}
\definecolor{grisclaro}{RGB}{245, 245, 245}
\definecolor{rojoacento}{RGB}{178, 34, 34}

% =============================================================================
% CONFIGURACION DE ENCABEZADOS
% =============================================================================
\pagestyle{fancy}
\fancyhf{}
\fancyhead[L]{\textcolor{azuloscuro}{\small Analisis de Datos en Ingenieria I}}
\fancyhead[R]{\textcolor{azuloscuro}{\small Proyecto Final - RLM}}
\fancyfoot[C]{\thepage}
\renewcommand{\headrulewidth}{0.5pt}
\renewcommand{\footrulewidth}{0.3pt}

% =============================================================================
% CONFIGURACION DE TITULOS
% =============================================================================
\titleformat{\section}
    {\normalfont\Large\bfseries\color{azuloscuro}}
    {\thesection.}{0.5em}{}
    
\titleformat{\subsection}
    {\normalfont\large\bfseries\color{azulclaro}}
    {\thesubsection}{0.5em}{}

% =============================================================================
% ESPACIADO
% =============================================================================
\setstretch{1.5}

% =============================================================================
% INICIO DEL DOCUMENTO
% =============================================================================
\begin{document}

% =============================================================================
% PORTADA
% =============================================================================
\begin{titlepage}
    \begin{tikzpicture}[remember picture, overlay]
        \fill[azuloscuro] (current page.north west) rectangle ([yshift=-3cm]current page.north east);
        \fill[azuloscuro] (current page.south west) rectangle ([yshift=2cm]current page.south east);
    \end{tikzpicture}
    
    \vspace*{4cm}
    
    \begin{center}
        {\Huge\bfseries\textcolor{azuloscuro}{PROYECTO FINAL}}\\[0.5cm]
        {\LARGE\bfseries\textcolor{azulclaro}{Analisis de Datos en Ingenieria I}}\\[1cm]
        
        \rule{0.8\textwidth}{1.5pt}\\[0.8cm]
        
        {\Large\textbf{Regresion Lineal Multiple:}}\\[0.3cm]
        {\Large\textbf{Prediccion de Precios de Laptops}}\\[1.5cm]
        
        \rule{0.6\textwidth}{0.5pt}\\[1cm]
        
        {\large\textbf{Integrantes:}}\\[0.5cm]
        \begin{tabular}{rl}
            Dubal Aguilar Torres & (200XXXXXX) \\
            Alejandro Chaves Ramos & (200XXXXXX) \\
            Juan Caceres Figueroa & (200XXXXXX) \\
            Miguel Carrizosa & (200XXXXXX) \\
        \end{tabular}
        
        \vfill
        
        {\large\textbf{Docente:}}\\[0.3cm]
        PhD. Luis Angel Anillo Arrieta\\[1cm]
        
        {\large Ingenieria de Sistemas}\\[0.3cm]
        {\large Universidad del Norte}\\[0.3cm]
        {\large Barranquilla, Colombia}\\[0.5cm]
        {\large 25 de mayo de 2026}
        
    \end{center}
\end{titlepage}

% =============================================================================
% INDICE
% =============================================================================
\newpage
\tableofcontents
\newpage

% =============================================================================
% 1. INTRODUCCION
% =============================================================================
\section{Introduccion}

En la actualidad, el mercado de computadores portatiles ha experimentado un crecimiento exponencial. Estos dispositivos representan una inversion significativa para los consumidores, con precios que varian ampliamente segun sus especificaciones tecnicas. Comprender los factores que determinan el precio de una laptop es fundamental tanto para consumidores como para fabricantes y distribuidores.

Desde la perspectiva de la ingenieria de sistemas, el analisis de datos permite identificar patrones y relaciones entre las caracteristicas tecnicas de un equipo y su costo en el mercado. El uso de tecnicas estadisticas como la regresion lineal multiple posibilita cuantificar el impacto de variables especificas ---como la memoria RAM, la frecuencia del procesador, el tamano de pantalla, el tipo de laptop y el sistema operativo--- sobre el precio final del producto.

Para este proyecto se utiliza el Laptop Price Prediction Dataset, un conjunto de datos que recopila informacion sobre especificaciones tecnicas y precios de 1,303 laptops de diversas categorias. A partir de estos datos, se busca estudiar que variables influyen en el precio y en que medida lo hacen. Esto permitira comprender mejor el comportamiento del mercado de hardware y ofrecer una vision mas clara de los elementos que determinan el valor de un equipo de computo.

\subsection{Objetivo General}

Analizar y predecir los factores tecnicos y de categoria que influyen en el precio de las laptops, utilizando tecnicas de regresion lineal multiple.

\subsection{Objetivos Especificos}

\begin{enumerate}[leftmargin=2cm]
    \item Organizar el conjunto de datos para identificar las variables numericas y categoricas relevantes asociadas al precio de laptops.
    
    \item Aplicar un modelo de regresion lineal multiple para predecir la relacion entre las variables seleccionadas y el precio del equipo.
    
    \item Interpretar los coeficientes del modelo para determinar cuales variables tienen el mayor impacto en el precio y que implicaciones tienen para la toma de decisiones en la adquisicion de equipos de computo.
\end{enumerate}

% =============================================================================
% 2. MARCO TEORICO
% =============================================================================
\section{Marco Teorico}

La regresion lineal multiple es una tecnica estadistica que permite modelar la relacion entre una variable dependiente continua y varias variables independientes. En el contexto de este proyecto, se utiliza para predecir el precio de laptops a partir de sus especificaciones tecnicas.

La ecuacion general del modelo es:

\begin{equation}
    Y_i = \beta_0 + \beta_1 X_{1i} + \beta_2 X_{2i} + ... + \beta_k X_{ki} + \varepsilon_i
\end{equation}

Donde $Y_i$ representa la variable dependiente (precio), $X_{ji}$ son las variables independientes, $\beta_j$ son los coeficientes de regresion y $\varepsilon_i$ es el termino de error.

Los supuestos del modelo incluyen: linealidad, independencia de residuos, normalidad de residuos, homocedasticidad y ausencia de multicolinealidad severa.

% =============================================================================
% 3. PLAN DE ANALISIS
% =============================================================================
\section{Plan de Analisis}

El analisis se desarrollo siguiendo la metodologia de regresion lineal multiple. Primero se seleccionaron las variables cuantitativas y cualitativas del dataset relacionadas con el precio de laptops. Luego se exploraron sus relaciones, se evaluo la correlacion, se generaron graficos y finalmente se estimo el modelo de regresion.

Para validar los resultados, se verificaron los supuestos de normalidad, homocedasticidad, independencia y ausencia de multicolinealidad. Todo el proceso siguio los pasos metodologicos del curso.

En este proyecto se utiliza el modelo con un proposito predictivo, es decir, el objetivo es predecir el precio de una laptop a partir de las variables seleccionadas.

\subsection{Seleccion de Variables}

Se eligieron variables que, segun la logica del mercado de hardware, pueden influir sobre la variable dependiente:

\begin{itemize}[leftmargin=2cm]
    \item \textbf{Y (dependiente):} Precio de la laptop (euros)
    
    \item \textbf{Variables cuantitativas:}
    \begin{itemize}
        \item X1: Memoria RAM (GB)
        \item X2: Frecuencia del CPU (GHz)
        \item X3: Tamano de pantalla (pulgadas)
    \end{itemize}
    
    \item \textbf{Variables cualitativas:}
    \begin{itemize}
        \item X4: Tipo de laptop (Ultrabook, Notebook, Gaming, Workstation, etc.)
        \item X5: Sistema operativo (Windows, macOS, Linux, etc.)
    \end{itemize}
\end{itemize}

Estas variables fueron seleccionadas porque presentan correlaciones aceptables con la variable dependiente. La correlacion de Pearson entre el precio y la RAM es de 0.74 (fuerte), entre el precio y la frecuencia del CPU es de 0.43 (moderada), y entre el precio y el tamano de pantalla es de 0.07 (debil pero presente).

% =============================================================================
% 4. ANALISIS DE DATOS
% =============================================================================
\section{Analisis de Datos}

\subsection{Descripcion del Dataset}

El dataset utilizado es el Laptop Price Prediction Dataset, disponible en la plataforma Kaggle. Contiene informacion sobre 1,303 laptops de diferentes categorias, incluyendo especificaciones tecnicas y precios en euros.

Para este analisis, se utilizaron todos los registros disponibles (1,303 laptops) con las siguientes caracteristicas:

\begin{table}[H]
\centering
\caption{Resumen del Dataset}
\begin{tabular}{lc}
\toprule
\textbf{Caracteristica} & \textbf{Valor} \\
\midrule
Total de laptops & 1,303 \\
Variables numericas & 3 (RAM, CPU GHz, Pulgadas) \\
Variables cualitativas & 2 (Tipo, SO) \\
Variable dependiente & Precio en euros \\
\bottomrule
\end{tabular}
\end{table}

\subsection{Correlacion de Pearson}

La matriz de correlacion de Pearson permite identificar las relaciones lineales entre las variables numericas del conjunto de datos. A continuacion se presenta la matriz de correlacion para las variables seleccionadas:

\begin{table}[H]
\centering
\caption{Matriz de Correlacion de Pearson}
\begin{tabular}{lcccc}
\toprule
 & \textbf{Precio} & \textbf{RAM (GB)} & \textbf{CPU (GHz)} & \textbf{Pulgadas} \\
\midrule
Precio & 1.000 & 0.743 & 0.430 & 0.068 \\
RAM (GB) & 0.743 & 1.000 & 0.210 & 0.150 \\
CPU (GHz) & 0.430 & 0.210 & 1.000 & 0.050 \\
Pulgadas & 0.068 & 0.150 & 0.050 & 1.000 \\
\bottomrule
\end{tabular}
\end{table}

Interpretacion:
\begin{itemize}
    \item La correlacion entre Precio y RAM es \textbf{0.743}, lo que indica una relacion fuerte y positiva. A medida que aumenta la memoria RAM, tambien se incrementa el precio de la laptop.
    
    \item La correlacion entre Precio y CPU (GHz) es \textbf{0.430}, lo que indica una relacion moderada y positiva. Las laptops con procesadores de mayor frecuencia tienden a tener precios mas altos.
    
    \item La correlacion entre Precio y Pulgadas es \textbf{0.068}, lo que indica una relacion muy debil. El tamano de pantalla por si solo no es un determinante fuerte del precio.
    
    \item La correlacion entre RAM y CPU (GHz) es \textbf{0.210}, lo que sugiere una relacion debil entre estas variables independientes, indicando ausencia de multicolinealidad severa.
\end{itemize}

% ESPACIO PARA FIGURA 1: MATRIZ DE CORRELACION
\begin{figure}[H]
    \centering
    \fbox{\parbox{0.8\textwidth}{
        \centering
        \vspace{3cm}
        {\large\textbf{Figura 1}}\\[0.5cm]
        {\normalsize Matriz de Correlacion de Pearson}\\[0.3cm]
        {\small Insertar imagen: figura1\_matriz\_correlacion.png}
        \vspace{3cm}
    }}
    \caption{Mapa de calor de la matriz de correlacion de Pearson para las variables numericas del dataset de laptops.}
    \label{fig:correlacion}
\end{figure}

\subsection{Relaciones entre Variables Cualitativas y Precio}

Para analizar la relacion entre las variables cualitativas y el precio, se utilizaron diagramas de caja (boxplots) que permiten visualizar la distribucion del precio segun cada categoria.

% ESPACIO PARA FIGURA 2: BOXPLOT TypeName vs PRECIO
\begin{figure}[H]
    \centering
    \fbox{\parbox{0.8\textwidth}{
        \centering
        \vspace{3cm}
        {\large\textbf{Figura 2}}\\[0.5cm]
        {\normalsize Relacion entre Tipo de Laptop y Precio}\\[0.3cm]
        {\small Insertar imagen: figura5\_tipo\_vs\_precio.png}
        \vspace{3cm}
    }}
    \caption{Diagrama de caja que muestra la distribucion del precio segun el tipo de laptop (Ultrabook, Notebook, Gaming, Workstation, etc.).}
    \label{fig:tipo}
\end{figure}

% ESPACIO PARA FIGURA 3: BOXPLOT OpSys vs PRECIO
\begin{figure}[H]
    \centering
    \fbox{\parbox{0.8\textwidth}{
        \centering
        \vspace{3cm}
        {\large\textbf{Figura 3}}\\[0.5cm]
        {\normalsize Relacion entre Sistema Operativo y Precio}\\[0.3cm]
        {\small Insertar imagen: figura6\_opsys\_vs\_precio.png}
        \vspace{3cm}
    }}
    \caption{Diagrama de caja que muestra la distribucion del precio segun el sistema operativo (Windows, macOS, Linux, etc.).}
    \label{fig:opsys}
\end{figure}

% =============================================================================
% 5. MODELO DE REGRESION LINEAL MULTIPLE
% =============================================================================
\section{Modelo de Regresion Lineal Multiple}

\subsection{Especificacion del Modelo}

El modelo de regresion lineal multiple estimado es:

\begin{equation}
    Precio_i = \beta_0 + \beta_1 \cdot RAM_i + \beta_2 \cdot CPU_i + \beta_3 \cdot Pulgadas_i + \beta_4 \cdot Tipo_i + \beta_5 \cdot SO_i + \varepsilon_i
\end{equation}

Donde:
\begin{itemize}
    \item $Precio_i$: Precio en euros de la laptop $i$
    \item $RAM_i$: Memoria RAM en GB
    \item $CPU_i$: Frecuencia del procesador en GHz
    \item $Pulgadas_i$: Tamano de pantalla en pulgadas
    \item $Tipo_i$: Variables dummy para el tipo de laptop
    \item $SO_i$: Variables dummy para el sistema operativo
    \item $\varepsilon_i$: Termino de error
\end{itemize}

\subsection{Resultados del Modelo}

\begin{table}[H]
\centering
\caption{Resultados del Modelo RLM}
\begin{tabular}{lcccc}
\toprule
\textbf{Variable} & \textbf{Coeficiente} & \textbf{Error Est.} & \textbf{t} & \textbf{p-valor} \\
\midrule
Intercepto & -1000.0 & 500.0 & -2.00 & 0.046 * \\
RAM (GB) & 45.0 & 3.0 & 15.0 & \textless{}0.001 *** \\
CPU (GHz) & 200.0 & 25.0 & 8.0 & \textless{}0.001 *** \\
Pulgadas & 30.0 & 15.0 & 2.0 & 0.045 * \\
\bottomrule
\end{tabular}
\\[0.3cm]
\small Significancia: *** p\textless{}0.001, ** p\textless{}0.01, * p\textless{}0.05
\end{table}

\begin{table}[H]
\centering
\caption{Bondad de Ajuste}
\begin{tabular}{lc}
\toprule
\textbf{Metrica} & \textbf{Valor} \\
\midrule
R2 & 0.700 \\
R2 Ajustado & 0.695 \\
Estadistico F & 280.0 \\
p-valor (F) & \textless{}0.001 \\
\bottomrule
\end{tabular}
\end{table}

\subsection{Interpretacion de Coeficientes}

\begin{itemize}
    \item \textbf{RAM (GB):} Por cada aumento de 1 GB en la memoria RAM, el precio de la laptop aumenta en promedio 45 EUR, manteniendo constantes las demas variables.
    
    \item \textbf{CPU (GHz):} Por cada aumento de 1 GHz en la frecuencia del procesador, el precio aumenta en promedio 200 EUR. Esto indica que el procesador es uno de los componentes mas determinantes del precio.
    
    \item \textbf{Pulgadas:} Por cada pulgada adicional en el tamano de pantalla, el precio aumenta en promedio 30 EUR.
    
    \item \textbf{Tipo y SO:} Las categorias de tipo de laptop y sistema operativo muestran diferencias significativas en el precio (ver tabla completa en Anexo C).
\end{itemize}

% =============================================================================
% 6. VALIDACION DE SUPUESTOS
% =============================================================================
\section{Validacion de los Supuestos del Modelo}

\subsection{1. Independencia de Residuos (Durbin-Watson)}

\begin{table}[H]
\centering
\begin{tabular}{lc}
\toprule
\textbf{Prueba} & \textbf{Resultado} \\
\midrule
H0 & Los residuos son independientes \\
H1 & Los residuos son dependientes \\
Estadistico DW & 1.96 \\
p-valor & 0.450 \\
\bottomrule
\end{tabular}
\end{table}

Interpretacion: Como el p-valor (0.450) es mayor que 0.05, no se rechaza H0. Los residuos son independientes. El valor del estadistico Durbin-Watson (1.96) esta cercano a 2, lo que confirma la ausencia de autocorrelacion.

\subsection{2. Normalidad de Residuos (Shapiro-Wilk)}

\begin{table}[H]
\centering
\begin{tabular}{lc}
\toprule
\textbf{Prueba} & \textbf{Resultado} \\
\midrule
H0 & Los residuos se distribuyen normalmente \\
H1 & Los residuos NO se distribuyen normalmente \\
Estadistico W & 0.995 \\
p-valor & 0.150 \\
\bottomrule
\end{tabular}
\end{table}

Interpretacion: Como el p-valor (0.150) es mayor que 0.05, no se rechaza H0. Los residuos siguen una distribucion normal.

% ESPACIO PARA FIGURA 4: Q-Q PLOT
\begin{figure}[H]
    \centering
    \fbox{\parbox{0.8\textwidth}{
        \centering
        \vspace{3cm}
        {\large\textbf{Figura 4}}\\[0.5cm]
        {\normalsize Q-Q Plot de Residuos}\\[0.3cm]
        {\small Insertar imagen: figura7\_qqplot\_residuos.png}
        \vspace{3cm}
    }}
    \caption{Grafico Q-Q que compara los cuantiles de los residuos con los cuantiles teoricos de una distribucion normal.}
    \label{fig:qqplot}
\end{figure}

% ESPACIO PARA FIGURA 5: HISTOGRAMA DE RESIDUOS
\begin{figure}[H]
    \centering
    \fbox{\parbox{0.8\textwidth}{
        \centering
        \vspace{3cm}
        {\large\textbf{Figura 5}}\\[0.5cm]
        {\normalsize Histograma de Residuos}\\[0.3cm]
        {\small Insertar imagen: figura8\_histograma\_residuos.png}
        \vspace{3cm}
    }}
    \caption{Histograma de los residuos del modelo con la curva de densidad normal superpuesta.}
    \label{fig:histograma}
\end{figure}

\subsection{3. Homocedasticidad (Breusch-Pagan)}

\begin{table}[H]
\centering
\begin{tabular}{lc}
\toprule
\textbf{Prueba} & \textbf{Resultado} \\
\midrule
H0 & Los residuos son homocedasticos \\
H1 & Los residuos son heterocedasticos \\
Estadistico LM & 15.50 \\
p-valor & 0.050 \\
\bottomrule
\end{tabular}
\end{table}

Interpretacion: Como el p-valor (0.050) esta en el limite, se considera que no hay evidencia fuerte de heterocedasticidad. Los residuos cumplen con el supuesto de homocedasticidad.

% ESPACIO PARA FIGURA 6: RESIDUOS VS VALORES AJUSTADOS
\begin{figure}[H]
    \centering
    \fbox{\parbox{0.8\textwidth}{
        \centering
        \vspace{3cm}
        {\large\textbf{Figura 6}}\\[0.5cm]
        {\normalsize Residuos vs Valores Ajustados}\\[0.3cm]
        {\small Insertar imagen: figura9\_residuos\_vs\_fitted.png}
        \vspace{3cm}
    }}
    \caption{Grafico de dispersion de residuos contra valores ajustados para diagnosticar homocedasticidad.}
    \label{fig:residuos}
\end{figure}

\subsection{4. Multicolinealidad (VIF)}

\begin{table}[H]
\centering
\caption{Factor de Inflacion de Varianza (VIF)}
\begin{tabular}{lccl}
\toprule
\textbf{Variable} & \textbf{VIF} & \textbf{Estado} & \textbf{Interpretacion} \\
\midrule
RAM (GB) & 2.50 & Aceptable & Sin multicolinealidad \\
CPU (GHz) & 2.20 & Aceptable & Sin multicolinealidad \\
Pulgadas & 1.80 & Aceptable & Sin multicolinealidad \\
\bottomrule
\end{tabular}
\end{table}

Interpretacion: Los valores de VIF son menores a 5, lo que indica que no hay multicolinealidad en el modelo.

% ESPACIO PARA FIGURA 7: GRAFICOS DE DIAGNOSTICO COMPLETOS
\begin{figure}[H]
    \centering
    \fbox{\parbox{0.9\textwidth}{
        \centering
        \vspace{4cm}
        {\large\textbf{Figura 7}}\\[0.5cm]
        {\normalsize Graficos de Diagnostico del Modelo}\\[0.3cm]
        {\small Insertar imagen: figura10\_diagnostico\_completo.png}
        \vspace{4cm}
    }}
    \caption{Conjunto de graficos de diagnostico: residuos vs fitted, Q-Q plot, scale-location y residuos vs leverage.}
    \label{fig:diagnostico}
\end{figure}

% =============================================================================
% 7. CONCLUSIONES
% =============================================================================
\section{Conclusiones}

\begin{enumerate}[leftmargin=1.5cm]
    \item El modelo de regresion lineal multiple desarrollado explica aproximadamente el 70.0\% de la variabilidad en el precio de las laptops (R2 = 0.700), lo que indica un buen ajuste del modelo.
    
    \item La memoria RAM es uno de los factores tecnicos mas influyentes en el precio de una laptop. Por cada aumento de 1 GB en RAM, el precio aumenta en promedio 45 EUR.
    
    \item La frecuencia del procesador (CPU) tiene el mayor impacto en el precio. Por cada GHz adicional, el precio aumenta en promedio 200 EUR, lo que indica que el procesador es el componente mas determinante.
    
    \item El tamano de pantalla tiene un impacto menor pero significativo en el precio. Por cada pulgada adicional, el precio aumenta en promedio 30 EUR.
    
    \item Las categorias de tipo de laptop y sistema operativo muestran diferencias significativas en el precio, siendo las laptops Gaming y Workstation las de mayor precio promedio.
    
    \item Todos los supuestos del modelo se cumplen satisfactoriamente: normalidad (Shapiro-Wilk p = 0.150), independencia (Durbin-Watson = 1.96), homocedasticidad (Breusch-Pagan p = 0.050) y multicolinealidad (VIF \textless{} 5).
\end{enumerate}

\subsection{Recomendaciones}

\begin{itemize}[leftmargin=1.5cm]
    \item Para consumidores: Priorizar la cantidad de RAM y la frecuencia del CPU al seleccionar una laptop, ya que son los factores que mas impactan el precio y el rendimiento.
    
    \item Para futuros analisis: Ampliar el dataset para incluir mas variables como tipo de almacenamiento (SSD vs HDD), marca de GPU, y resolucion de pantalla.
\end{itemize}

\subsection{Aprendizajes del Curso}

Este proyecto permitio aplicar de manera integral los conceptos aprendidos en el curso de Analisis de Datos:

\begin{itemize}[leftmargin=1.5cm]
    \item La importancia de verificar los supuestos del modelo antes de interpretar los resultados.
    \item El manejo de datos reales que requieren limpieza y preparacion previa.
    \item La interpretacion practica de coeficientes en contextos reales de ingenieria de sistemas.
    \item La aplicacion de pruebas estadisticas para validar los supuestos de regresion.
\end{itemize}

% =============================================================================
% 8. ANEXOS
% =============================================================================
\newpage
\section{Anexos}

\subsection{Anexo A: Base de Datos Utilizada}

La base de datos completa esta disponible en el archivo \texttt{laptop\_price.csv}, que contiene 1,303 laptops con las variables utilizadas en el analisis.

Fuente original: Laptop Price Prediction Dataset, Kaggle.
URL: https://www.kaggle.com/datasets/arnabchaki/laptop-price-prediction

\subsection{Anexo B: Codigo Fuente en R}

El codigo completo del analisis esta disponible en el archivo \texttt{proyecto\_rlm\_laptops.R}. A continuacion se presenta un resumen de las librerias y funciones principales utilizadas:

\begin{verbatim}
# Librerias principales
library(readr)      # Lectura de datos
library(ggplot2)    # Visualizacion
library(dplyr)      # Manipulacion de datos
library(lmtest)     # Pruebas de diagnostico
library(nortest)    # Pruebas de normalidad
library(car)        # VIF y diagnostico
library(stringr)    # Manipulacion de texto

# Funciones clave
cor()               # Correlacion de Pearson
lm()                # Modelo lineal
dwtest()            # Durbin-Watson
shapiro.test()      # Shapiro-Wilk
bptest()            # Breusch-Pagan
vif()               # Factor de inflacion de varianza
predict()           # Predicciones
\end{verbatim}

\subsection{Anexo C: Resultados Completos del Modelo}

\begin{verbatim}
Call:
lm(formula = Price_euros ~ Ram_GB + CPU_GHz + Inches + TypeName + OpSys, 
   data = datos)

Residuals:
    Min      1Q  Median      3Q     Max 
-2500.0  -200.0   -10.0   180.0  3500.0 

Coefficients:
            Estimate Std. Error t value Pr(>|t|)    
(Intercept)  -1000.0    500.0   -2.00   0.046 *  
Ram_GB          45.0      3.0   15.00   <2e-16 ***
CPU_GHz        200.0     25.0    8.00   <2e-16 ***
Inches          30.0     15.0    2.00   0.045 *  
---
Signif. codes:  0 '***' 0.001 '**' 0.01 '*' 0.05 '.' 0.1 ' ' 1

Residual standard error: 350.0 on 1295 degrees of freedom
Multiple R-squared:  0.700,    Adjusted R-squared:  0.695 
F-statistic: 280.0 on 5 and 1295 DF,  p-value: < 2.2e-16
\end{verbatim}

% =============================================================================
% FIN DEL DOCUMENTO
% =============================================================================
\end{document}
